% GENERATED FILE. Edit canonical lessons, then run npm run generate:books.
% canonical-source-hash: fnv1a64:ce42f7ec79dc3633
% canonical-lessons: GU-C510-soni, GU-C510-ogaalvu, GU-C510-tabu, GU-C510-dhaalvu, GU-C510-ghadavvu

\chapter{At the Goldsmith's}
\label{ch:gu-at-the-goldsmith-s}
\hlchaptermodality{510}

\begin{chapteropening}
\textbf{By the end of this chapter:} \emph{I can talk about what a goldsmith melts, casts and makes.}

Complete the last lesson of chapter 510: I can talk about what a goldsmith melts, casts and makes.
\end{chapteropening}

\section[\texorpdfstring{\gu{સોની} (sonī)}{સોની (sonī)}]{\gu{સોની} (sonī) — a goldsmith}

\label{lesson:GU-C510-soni}



\begin{quote}
Before the new one: say the Gujarati for rust, then the Gujarati for to rust.
\end{quote}

\subsection*{You'll want to know: \gu{સોની}}
\textbf{\gu{સોની}} — \emph{sonī} — \textquotedblleft{}a goldsmith\textquotedblright{}.

\begin{quote}
\gu{અમારી} \gu{ગલીમાં} \gu{એક} \gu{સોનીની} \gu{દુકાન} \gu{છે}.
\end{quote}

\begin{quote}
\emph{amārī galīmā̃ ek sonīnī dukān che.}
\end{quote}

\begin{quote}
There is a goldsmith's shop in our lane.
\end{quote}

A masculine noun.

\begin{cousinweb}
\textbf{\gu{સોની}} (\emph{sonī}) is built on the Gujarati word for gold, which goes back to Sanskrit \emph{suvárṇa}, \textquotedblleft{}gold\textquotedblright{} (CDIAL 13519). CDIAL also lists the older \emph{sonār}, \textquotedblleft{}a goldsmith\textquotedblright{}, from Sanskrit \emph{suvarṇakāra} (CDIAL 13520).
\end{cousinweb}

\begin{grammarlens}[title={What you've built}]
One more word for describing what you lived and felt.
\end{grammarlens}

\subsection*{Your turn}
\begin{itemize}
\raggedright
  \item \emph{Say it:} \emph{sonī}
  \item \emph{Say it:} \emph{sonī}, once more
  \item \emph{Say it:} \emph{amārī galīmā̃ ek sonīnī dukān che.} — There is a goldsmith's shop in our lane.
  \item \emph{Recall:} say the Gujarati for rust, then the Gujarati for to rust, then say \emph{sonī} again
\end{itemize}

\begin{tcolorbox}[breakable,colback=teal!4,colframe=teal!35!black,title={Before you move on}]
What does \gu{સોની} mean? (\textquotedblleft{}A goldsmith\textquotedblright{}.) Say it once more.
\end{tcolorbox}
\section[\texorpdfstring{\gu{ઓગાળવું} (ogāḷvũ)}{ઓગાળવું (ogāḷvũ)}]{\gu{ઓગાળવું} (ogāḷvũ) — to melt down (metal)}

\label{lesson:GU-C510-ogaalvu}



\begin{quote}
Before the new one: say the Gujarati for to rust, then the Gujarati for a goldsmith.
\end{quote}

\subsection*{You'll want to know: \gu{ઓગાળવું}}
\textbf{\gu{ઓગાળવું}} — \emph{ogāḷvũ} — \textquotedblleft{}to melt down (metal)\textquotedblright{}.

\begin{quote}
\gu{સોનીએ} \gu{જૂની} \gu{વીંટી} \gu{ઓગાળી}.
\end{quote}

\begin{quote}
\emph{sonīe jūnī vī̃ṭī ogāḷī.}
\end{quote}

\begin{quote}
The goldsmith melted down the old ring.
\end{quote}

Past, with \gu{તેણે} (he or she …): \gu{તેણે} \gu{વીંટી} \gu{ઓગાળી}, \emph{teṇe vī̃ṭī ogāḷī}.

\begin{cousinweb}
\textbf{\gu{ઓગાળવું}} (\emph{ogāḷvũ}) goes back to Sanskrit \emph{avagalati}, \textquotedblleft{}falls down, slips down\textquotedblright{}; CDIAL lists the Gujarati verb for making something melt beside the plain verb for oozing and melting (CDIAL 742). The past agrees with what is melted: \textbf{\gu{વીંટી} \gu{ઓગાળી}}.
\end{cousinweb}

\begin{grammarlens}[title={What you've built}]
One more word for describing what you lived and felt.
\end{grammarlens}

\subsection*{Your turn}
\begin{itemize}
\raggedright
  \item \emph{Say it:} \emph{ogāḷvũ}
  \item \emph{Say it:} \emph{ogāḷvũ}, once more
  \item \emph{Say it:} \emph{sonīe jūnī vī̃ṭī ogāḷī.} — The goldsmith melted down the old ring.
  \item \emph{Recall:} say the Gujarati for to rust, then the Gujarati for a goldsmith, then say \emph{ogāḷvũ} again
\end{itemize}

\begin{tcolorbox}[breakable,colback=teal!4,colframe=teal!35!black,title={Before you move on}]
What does \gu{ઓગાળવું} mean? (\textquotedblleft{}To melt down (metal)\textquotedblright{}.) Say it once more.
\end{tcolorbox}
\section[\texorpdfstring{\gu{તાંબું} (tā̃bũ)}{તાંબું (tā̃bũ)}]{\gu{તાંબું} (tā̃bũ) — copper}

\label{lesson:GU-C510-tabu}



\begin{quote}
Before the new one: say the Gujarati for a goldsmith, then the Gujarati for to melt down.
\end{quote}

\subsection*{You'll want to know: \gu{તાંબું}}
\textbf{\gu{તાંબું}} — \emph{tā̃bũ} — \textquotedblleft{}copper\textquotedblright{}.

\begin{quote}
\gu{દાદીનું} \gu{જૂનું} \gu{વાસણ} \gu{તાંબાનું} \gu{છે}.
\end{quote}

\begin{quote}
\emph{dādīnũ jūnũ vāsaṇ tā̃bānũ che.}
\end{quote}

\begin{quote}
Grandmother's old pot is made of copper.
\end{quote}

A neuter noun.

\begin{cousinweb}
\textbf{\gu{તાંબું}} (\emph{tā̃bũ}) is a tadbhava, inherited from Sanskrit \emph{tāmrá}, \textquotedblleft{}coppery red; copper\textquotedblright{}, and CDIAL lists the Gujarati word (CDIAL 5779).
\end{cousinweb}

\begin{grammarlens}[title={What you've built}]
One more word for describing what you lived and felt.
\end{grammarlens}

\subsection*{Your turn}
\begin{itemize}
\raggedright
  \item \emph{Say it:} \emph{tā̃bũ}
  \item \emph{Say it:} \emph{tā̃bũ}, once more
  \item \emph{Say it:} \emph{dādīnũ jūnũ vāsaṇ tā̃bānũ che.} — Grandmother's old pot is made of copper.
  \item \emph{Recall:} say the Gujarati for a goldsmith, then the Gujarati for to melt down, then say \emph{tā̃bũ} again
\end{itemize}

\begin{tcolorbox}[breakable,colback=teal!4,colframe=teal!35!black,title={Before you move on}]
What does \gu{તાંબું} mean? (\textquotedblleft{}Copper\textquotedblright{}.) Say it once more.
\end{tcolorbox}
\section[\texorpdfstring{\gu{ઢાળવું} (ḍhāḷvũ)}{ઢાળવું (ḍhāḷvũ)}]{\gu{ઢાળવું} (ḍhāḷvũ) — to cast (metal in a mould)}

\label{lesson:GU-C510-dhaalvu}



\begin{quote}
Before the new one: say the Gujarati for to melt down, then the Gujarati for copper.
\end{quote}

\subsection*{You'll want to know: \gu{ઢાળવું}}
\textbf{\gu{ઢાળવું}} — \emph{ḍhāḷvũ} — \textquotedblleft{}to cast (metal in a mould)\textquotedblright{}.

\begin{quote}
\gu{સોનીએ} \gu{ગરમ} \gu{તાંબું} \gu{એક} \gu{નાના} \gu{વાસણમાં} \gu{ઢાળ્યું}.
\end{quote}

\begin{quote}
\emph{sonīe garam tā̃bũ ek nānā vāsaṇmā̃ ḍhāḷyũ.}
\end{quote}

\begin{quote}
The goldsmith poured the hot copper into a small pot.
\end{quote}

Past, with \gu{તેણે} (he or she …): \gu{તેણે} \gu{તાંબું} \gu{ઢાળ્યું}, \emph{teṇe tā̃bũ ḍhāḷyũ}.

\begin{cousinweb}
\textbf{\gu{ઢાળવું}} (\emph{ḍhāḷvũ}) comes from \emph{*ḍhalati}, a form CDIAL reconstructs, and CDIAL lists the Gujarati verb as \textquotedblleft{}to pour\textquotedblright{} (CDIAL 5581): metal poured into a mould is cast.
\end{cousinweb}

\begin{grammarlens}[title={What you've built}]
One more word for describing what you lived and felt.
\end{grammarlens}

\subsection*{Your turn}
\begin{itemize}
\raggedright
  \item \emph{Say it:} \emph{ḍhāḷvũ}
  \item \emph{Say it:} \emph{ḍhāḷvũ}, once more
  \item \emph{Say it:} \emph{sonīe garam tā̃bũ ek nānā vāsaṇmā̃ ḍhāḷyũ.} — The goldsmith poured the hot copper into a small pot.
  \item \emph{Recall:} say the Gujarati for to melt down, then the Gujarati for copper, then say \emph{ḍhāḷvũ} again
\end{itemize}

\begin{tcolorbox}[breakable,colback=teal!4,colframe=teal!35!black,title={Before you move on}]
What does \gu{ઢાળવું} mean? (\textquotedblleft{}To cast (metal in a mould)\textquotedblright{}.) Say it once more.
\end{tcolorbox}
\section[\texorpdfstring{\gu{ઘડાવવું} (ghaḍāvvũ)}{ઘડાવવું (ghaḍāvvũ)}]{\gu{ઘડાવવું} (ghaḍāvvũ) — to have (something) made}

\label{lesson:GU-C510-ghadavvu}



\begin{quote}
Before the new one: say the Gujarati for copper, then the Gujarati for to cast.
\end{quote}

\subsection*{You'll want to know: \gu{ઘડાવવું}}
\textbf{\gu{ઘડાવવું}} — \emph{ghaḍāvvũ} — \textquotedblleft{}to have (something) made\textquotedblright{}.

\begin{quote}
\gu{મીરાએ} \gu{સોની} \gu{પાસે} \gu{એક} \gu{વીંટી} \gu{ઘડાવી}.
\end{quote}

\begin{quote}
\emph{Mīrāe sonī pāse ek vī̃ṭī ghaḍāvī.}
\end{quote}

\begin{quote}
Mira had a ring made by the goldsmith.
\end{quote}

Past, with \gu{તેણે} (he or she …): \gu{તેણે} \gu{વીંટી} \gu{ઘડાવી}, \emph{teṇe vī̃ṭī ghaḍāvī}.

\begin{cousinweb}
\textbf{\gu{ઘડાવવું}} (\emph{ghaḍāvvũ}) is built from \textbf{\gu{ઘડવું}}, \textquotedblleft{}to shape\textquotedblright{}, with the causative \emph{-āv-}: to have something shaped by another person, who takes \textbf{\gu{પાસે}}. The family goes back to Sanskrit \emph{ghaṭ}, whose causative \emph{ghaṭayati} Monier-Williams gives as \textquotedblleft{}to join, to put together\textquotedblright{} (Monier-Williams).
\end{cousinweb}

\begin{grammarlens}[title={What you've built}]
One more word for describing what you lived and felt.
\end{grammarlens}

\subsection*{Your turn}
\begin{itemize}
\raggedright
  \item \emph{Say it:} \emph{ghaḍāvvũ}
  \item \emph{Say it:} \emph{ghaḍāvvũ}, once more
  \item \emph{Say it:} \emph{Mīrāe sonī pāse ek vī̃ṭī ghaḍāvī.} — Mira had a ring made by the goldsmith.
  \item \emph{Recall:} say the Gujarati for copper, then the Gujarati for to cast, then say \emph{ghaḍāvvũ} again
\end{itemize}

\begin{tcolorbox}[breakable,colback=teal!4,colframe=teal!35!black,title={Before you move on}]
What does \gu{ઘડાવવું} mean? (\textquotedblleft{}To have (something) made\textquotedblright{}.) Say it once more.
\end{tcolorbox}
